\movedhead{Comparison to human baselines}
\label{sec:humanbaselines}\label{app:humanbaselines}
Two prior measurements of the disposition effect in language models should be
read before the comparison below, because neither reproduces the human
direction. \citet{malberg2025comprehensive}, testing thirty biases across
twenty models, report the disposition effect as one of only two with a strong
negative average, meaning their models preferred change to the status quo.
\citet{yee2026calibrating}, across four models and 24{,}000
agent-scenario pairs, report disposition effects near zero relative to human
benchmarks. A model that is not reluctant to realise a loss is therefore the
established finding on the model side, measured on two instruments that are
not this one, and H1 was registered in the human direction in spite of it.

The human benchmark for the disposition family is
\citet{odean1998losses}, who scores 10,000 accounts at a nationwide discount
brokerage from 1987 to 1993 by the proportion of gains realised against the
proportion of losses realised. Aggregated over the whole period his Table~I
reports PGR of 0.148 against PLR of 0.098, a difference of 0.050 with a
$t$-statistic of 35 over 13,883 realised gains, 79,658 paper gains, 11,930
realised losses and 110,348 paper losses. Investors realised their winners
about half again as often as their losers.

That is not the quantity this battery measures, and the two should not be
placed in the same column. Odean's denominator is the set of positions
available to be sold, which only a portfolio can supply; a prompt cannot.
What transfers between the two instruments is the sign. A human investor is
more willing to sell a position held at a gain than one held at a loss, so
the disposition family here predicts a lower sell tendency in the loss arm
than in the gain arm, which is the negative direction registered for H1 in
Table~\ref{tab:hypotheses}. Whether the model shares that sign is a question
the matched pairs can answer even though the magnitudes are not comparable.

The models do not share that sign. H1 predicted a negative effect, a lower
sell tendency in the loss arm, and the measured effect is positive in eight
of the nine model by effort cells on both measures. The exception is Opus at
\texttt{low} effort, discussed below. Aggregating the disposition family over
all nine cells, which the analysis plan does not register as a quantity and
which is reported here as a secondary aggregation rather than a confirmatory
test, the effect is $+0.239$ $[+0.148, +0.331]$ on the categorical measure
over 108 pairs, 47 of them positive against 11 negative, and $+13.376$
$[+8.844, +18.104]$ on the continuous measure, 64 positive against 16
negative. Per model, the aggregate is $+0.367$ $[+0.206, +0.533]$ in Haiku,
$+0.328$ $[+0.178, +0.500]$ in Sonnet, and $+0.022$ $[-0.106, +0.144]$ in
Opus.

Read against Odean's sign, this is a reversal rather than an absence. A
position held below its cost basis draws more selling from these models, not
less, and in the two smaller models the reversal is large and its interval is
far from zero. H1 is recorded as a failed directional prediction. The
pre-registered sign in \texttt{analysis/directions.py} was not edited after
the data arrived, and the aggregation above is reported with the sign the
design predicted still attached to it, so the failure is visible in the
table rather than absorbed into it.

Two independent instruments already pointed this way.
\citet{malberg2025comprehensive} score the disposition effect as one of only
two biases out of thirty with a strong negative average, which on their
coding means a preference for change, and \citet{yee2026calibrating} find
disposition effects near zero against human benchmarks. The present
measurement is the third, on a third instrument, and it is the first to
attach an interval to the reversal in an advisory setting under decision
pressure.
